% §6.6 Bibliographic Notes
%
% Cross-reference resolution status:
%
% ALREADY RESOLVED (within chapter 6, drafted in Tasks 2-6 of Plan 6):
%   \Cref{sec:union-of-sets}                    -- chapter 6 §6.1
%   \Cref{sec:other-set-ops}                    -- chapter 6 §6.2
%   \Cref{sec:cancellation}                     -- chapter 6 §6.3
%   \Cref{sec:perfect-filter-line}              -- chapter 6 §6.4
%   \Cref{prop:complement-swap}                 -- chapter 6 §6.2
%
% ALREADY RESOLVED (chapter 3, drafted in Plan 3):
%   \Cref{sec:or-not}                           -- chapter 3 §3.2
%   \Cref{sec:composition-theorem}              -- chapter 3 §3.3
%
% ALREADY RESOLVED (chapter 5, drafted in Plan 5):
%   \Cref{sec:predicate-algebra}                -- chapter 5 §5.3
%   \Cref{sec:one-sided-two-sided}              -- chapter 5 §5.4
%
% ALREADY RESOLVED (Plan 0 chapter stubs):
%   \Cref{ch:perfect-hash-filters}              -- chapter 15 stub label
%
% Citations resolved by new references.bib entries added in this task:
%   \cite{towell-bernoulli-composition}         -- Towell Bernoulli composition manuscript
%
% Citations resolved by prior references.bib entries:
%   \cite{bloomfilter}                          -- already in references.bib
%   \cite{mitzenmacher-upfal}                   -- already in references.bib (Plan 3)
%   \cite{broder-mitzenmacher-2004}             -- already in references.bib (Plan 5)
%   \cite{carter-wegman-1979}                   -- already in references.bib (Plan 5)
%   \cite{towell-bernoulli-sets}                -- already in references.bib (Plan 0)
%   \cite{towell-bernoulli-hash-function}       -- already in references.bib (Plan 5)

\section*{Bibliographic Notes}
\addcontentsline{toc}{section}{Bibliographic Notes}
\label{sec:bib-notes-ch6}

\paragraph{The composition framework.}
The rate-propagation formulas of \Cref{sec:union-of-sets,sec:other-set-ops}
are the set-level instances of a more general framework developed in
\textcite{towell-bernoulli-composition}. That manuscript proves the
rate-propagation laws for the full algebra of approximate sets---union,
intersection, complement, difference, and the symmetric difference---under
independence of the constituent $\Bernoulli{\BoolType}$ outputs, and
catalogues which compound operations preserve and which destroy
one-sidedness. The chapter's contribution is to derive each rule
constructively from the predicate algebra of \Cref{sec:predicate-algebra}
and the AND/OR/NOT rules of \Cref{sec:or-not}, rather than stating the
formulas as a fait accompli. The composition manuscript treats the rules
at the level of joint $(\fprate, \fnrate)$ pairs; chapter~6 unpacks the
proofs and connects each rule back to its Boolean-algebra ancestor in
chapter~3.

\paragraph{Bloom filter union and intersection.}
The well-known fact that two Bloom filters of the same length and the same
hash family admit a bitwise union (logical OR of the bit arrays) and a
bitwise intersection (logical AND) is folklore by now, but the analysis
of what the resulting filter \emph{computes}---and at what error rate---is
treated systematically in \textcite{mitzenmacher-upfal}, Chapter~5,
and surveyed across applications in
\textcite{broder-mitzenmacher-2004}. The union case is clean: the bitwise
OR of two filters for sets $A$ and $B$ is exactly a filter for $A \cup B$,
with $\fnrate = 0$ preserved and a slightly elevated $\fprate$ that
chapter~6's \Cref{prop:union-fpr-fnr} accounts for. The intersection case
is famously subtle: the bitwise AND of two filters does \emph{not} in
general produce a filter for $A \cap B$, because the bit pattern of an AND
filter admits ghost elements not in either input. Chapter~6 sidesteps the
bitwise-AND trap by treating intersection at the predicate level instead,
which gives the correct rate without committing to a specific
representation. Readers who want the bitwise variant and its diagnostic
issues should consult the survey directly.

\paragraph{De Morgan under noise.}
The complement-swap result of \Cref{prop:complement-swap}---that
complementing an approximate set exchanges its FPR and FNR---is the set-level
shadow of the Boolean-algebra rule worked out in \Cref{sec:or-not}: NOT
swaps the two error rates of a noisy bit. Classical De Morgan reads
$\complement (A \cup B) = \complement A \cap \complement B$, an identity
between sets. Under noise the identity at the set level still holds, but
the error rates of the two sides traverse the rate plane along different
paths: one side touches the FNR-zero axis if both inputs are FNR-zero,
the other side touches the FPR-zero axis after the swap. The
``different paths, same destination'' phenomenon is what makes the
perfect-filter-line classification of \Cref{sec:perfect-filter-line}
informative: it tells you which compound operations stay on a one-sided
axis and which mixed operations are forced into the interior. This
chapter's treatment is the set-algebra rephrasing of the noisy-bit
De~Morgan worked through in chapter~3.

\paragraph{Perfect-filter literature and the bridge forward.}
The perfect-filter line of \Cref{sec:perfect-filter-line}---the locus of
one-sided sets in the rate plane (\Cref{sec:one-sided-two-sided})---is
where chapter~6 lands its main classification result. The line itself is
a chapter~5 construct, but the question of \emph{which} compound operations
preserve membership on the line is the chapter~6 contribution. The
constructive side of the story, how to actually build a set whose rates
sit exactly on the FPR-zero axis, has its roots in the perfect-hashing
literature opened by \textcite{carter-wegman-1979} and is taken up in full
in chapter~15 (\Cref{ch:perfect-hash-filters}), which builds on
\textcite{towell-bernoulli-hash-function} for the space-optimal construction.
The Bloom filter of \textcite{bloomfilter} occupies the opposite axis;
it has been the canonical one-sided approximate set since 1970, and its
algebraic behaviour under union and intersection is exactly what this
chapter's closure analysis formalizes. Chapter~6 sits between the
predicate-algebra reading of chapter~5 and the constructive perfect-filter
machinery of chapter~15: the algebra of operations, divorced from any
specific representation.
